%==============================================================================
% Chapitre 13 - Rotation terrestre et effets géophysiques
%==============================================================================

\section{Introduction}

Dans la plupart des situations courantes, la Terre peut être considérée comme
immobile.

Cette approximation est suffisante pour de nombreuses applications
balistiques.

Cependant, lorsque les distances augmentent et que la précision recherchée
devient importante, plusieurs phénomènes liés à la rotation terrestre
apparaissent.

Parmi les plus connus figurent :

\begin{itemize}
    \item l'effet Coriolis ;
    \item l'effet Eötvös ;
    \item la courbure terrestre ;
    \item les variations locales de la gravité.
\end{itemize}

Ces phénomènes sont généralement négligeables à courte distance mais peuvent
devenir mesurables lors des tirs à longue portée.

\begin{aretenir}

La rotation terrestre influence la trajectoire des projectiles.

Son effet est généralement faible à courte distance mais devient mesurable à
longue portée.

\end{aretenir}

\section{La Terre est en rotation}

La Terre effectue une rotation complète sur elle-même en environ :

\begin{equation}
23~h~56~min~4~s
\end{equation}

Cette durée correspond au jour sidéral.

La vitesse de rotation dépend de la latitude.

Elle est maximale à l'équateur et nulle aux pôles.

\begin{exemple}

À l'équateur, un point situé à la surface terrestre se déplace à environ :

\begin{equation}
465~\mathrm{m/s}
\end{equation}

vers l'est.

\end{exemple}

\section{Référentiel terrestre et référentiel inertiel}

Comme indiqué au chapitre~\ref{chap:referentiels}, la Terre n'est pas un
référentiel parfaitement inertiel.

Sa rotation introduit des accélérations qui peuvent produire des effets
apparents sur les trajectoires.

Pour les calculs balistiques de haute précision, ces effets doivent parfois
être pris en compte.

\section{Principe général}

Lorsqu'un projectile quitte le canon, il conserve la vitesse qu'il possède au
moment du tir.

Pendant son vol, la Terre continue de tourner sous sa trajectoire.

Cette situation conduit à l'apparition de plusieurs phénomènes observables
dans le référentiel terrestre.

\begin{exemple}

Un projectile tiré vers le nord depuis l'hémisphère nord conserve initialement
la vitesse de rotation correspondant à sa latitude de départ.

Or les régions situées plus au nord possèdent une vitesse tangentielle
légèrement plus faible.

Cette différence contribue à la déviation observée.

\end{exemple}

\section{Effet Coriolis}

L'effet Coriolis est une conséquence directe de la rotation terrestre.

Dans l'hémisphère nord, il tend à dévier les trajectoires :

\begin{itemize}
    \item vers la droite du mouvement ;
\end{itemize}

alors que dans l'hémisphère sud il tend à les dévier :

\begin{itemize}
    \item vers la gauche du mouvement.
\end{itemize}

\begin{aretenir}

Dans l'hémisphère nord, l'effet Coriolis produit généralement une déviation
vers la droite.

\end{aretenir}

\section{Compréhension intuitive}

Une façon simple de comprendre l'effet Coriolis consiste à considérer les
différences de vitesse tangentielle entre les latitudes.

Les points situés près de l'équateur se déplacent plus rapidement vers l'est
que ceux situés près des pôles.

Un projectile tiré vers le nord conserve initialement la vitesse tangentielle
de son point de départ.

Il se retrouve alors légèrement décalé vers l'est par rapport au terrain
survolé.

Cette explication est utile mais ne décrit pas tous les cas possibles.

L'approche complète repose sur la mécanique dans un référentiel en rotation.

\section{Expression mathématique}

\begin{rigueur}

L'accélération de Coriolis s'écrit :

\begin{equation}
\vec a_C =
-2\,\vec\Omega \times \vec v
\label{eq:coriolis}
\end{equation}

où :

\begin{itemize}
    \item $\vec a_C$ est l'accélération de Coriolis ;
    \item $\vec\Omega$ est le vecteur vitesse angulaire terrestre ;
    \item $\vec v$ est la vitesse du projectile.
\end{itemize}

\end{rigueur}

Cette formulation sera réutilisée dans les chapitres consacrés aux modèles
balistiques avancés.

\section{Effet Eötvös}

L'effet Eötvös est moins connu que l'effet Coriolis.

Il apparaît lorsqu'un projectile possède une composante de vitesse vers
l'est ou vers l'ouest.

Lors d'un tir vers l'est :

\begin{itemize}
    \item la vitesse de rotation apparente augmente ;
    \item le poids apparent diminue légèrement.
\end{itemize}

Lors d'un tir vers l'ouest :

\begin{itemize}
    \item la vitesse de rotation apparente diminue ;
    \item le poids apparent augmente légèrement.
\end{itemize}

\begin{exemple}

Deux projectiles identiques tirés à la même vitesse :

\begin{itemize}
    \item l'un vers l'est ;
    \item l'autre vers l'ouest ;
\end{itemize}

ne suivront pas exactement la même trajectoire.

\end{exemple}

\section{Courbure terrestre}

À courte distance, la surface terrestre peut être considérée comme plane.

Cette approximation devient progressivement moins valide lorsque la distance
augmente.

La ligne de visée, la trajectoire du projectile et la surface terrestre ne
partagent alors plus la même géométrie.

La courbure terrestre intervient notamment :

\begin{itemize}
    \item dans les tirs à très longue distance ;
    \item dans l'artillerie ;
    \item dans les simulations de grande échelle.
\end{itemize}

\section{Variations de la gravité}

La valeur usuelle :

\begin{equation}
g \approx 9,81~\mathrm{m/s^2}
\end{equation}

n'est qu'une moyenne.

La gravité varie légèrement :

\begin{itemize}
    \item avec la latitude ;
    \item avec l'altitude ;
    \item avec la forme de la Terre.
\end{itemize}

Pour les armes légères, ces variations sont généralement négligeables.

\section{Importance relative des effets}

Tous les phénomènes étudiés dans ce chapitre n'ont pas la même importance.

Pour les armes légères, l'ordre d'importance est souvent :

\begin{enumerate}
    \item gravité ;
    \item traînée ;
    \item vent ;
    \item stabilité gyroscopique ;
    \item effet Coriolis ;
    \item effet Eötvös ;
    \item variations locales de la gravité.
\end{enumerate}

Cette hiérarchie dépend naturellement :

\begin{itemize}
    \item de la distance ;
    \item de la précision recherchée ;
    \item du type de projectile.
\end{itemize}

\section{Application aux simulations}

Les modèles balistiques simplifiés ignorent généralement les effets liés à la
rotation terrestre.

Les modèles avancés peuvent intégrer :

\begin{itemize}
    \item l'effet Coriolis ;
    \item l'effet Eötvös ;
    \item la courbure terrestre ;
    \item des modèles gravitationnels plus complets.
\end{itemize}

Ces effets deviennent particulièrement intéressants dans les simulations de
tirs à longue portée.

\begin{simulation}

Pour une simulation d'armes légères à quelques centaines de mètres, les effets
liés à la rotation terrestre sont souvent négligeables devant ceux du vent.

À plusieurs kilomètres, ils peuvent devenir mesurables et doivent parfois être
pris en compte.

\end{simulation}

\section{À retenir}

\begin{aretenir}

\begin{itemize}
    \item La Terre n'est pas un référentiel parfaitement inertiel.
    \item La rotation terrestre engendre plusieurs effets mesurables.
    \item L'effet Coriolis provoque une déviation latérale des trajectoires.
    \item L'effet Eötvös modifie légèrement le poids apparent.
    \item La courbure terrestre devient importante aux grandes distances.
    \item Ces effets sont généralement secondaires par rapport au vent pour
    les armes légères.
\end{itemize}

\end{aretenir}
